\chapter{Testing}

Testing was conducted to validate the functionality, accuracy, and performance of ProblySUS. A layered testing methodology was adopted — beginning with individual module unit testing, followed by integration testing of the full analysis pipeline, and concluding with performance and usability evaluations under real-world URL scanning scenarios.

\section{Unit Testing}

Unit testing was carried out to confirm that each analysis module operated correctly in isolation before being integrated into the full pipeline. The primary objective was to verify that individual components produced accurate, expected outputs across a range of valid and edge-case inputs.

A custom \texttt{test\_scorer.py} script was written to unit-test the scoring engine with synthetic \texttt{check\_results} dictionaries — verifying that each signal category produced the correct penalty, compound signals triggered their amplification bonus, and blacklist hits returned the hard override of score 100 / Fraudulent regardless of other signals.

Individual modules such as \texttt{pattern\_checker.py} and \texttt{validator.py} were tested with crafted URLs covering IP-based hosts, abused TLDs, excessive hyphens, missing schemes, and invalid syntax.

\begin{table}[h]
\centering
\begin{tabular}{|p{0.22\textwidth}|p{0.60\textwidth}|p{0.10\textwidth}|}
\toprule
Module & Test Case Description & Result \\
\midrule
validator.py & Accepts valid URLs, rejects malformed inputs, normalises missing schemes & Pass \\
pattern\_checker.py & Detects IP hosts, abused TLDs, hyphens, and phishing keywords correctly & Pass \\
whois\_checker.py & Returns correct domain age in days; handles WHOIS unavailable gracefully & Pass \\
dns\_checker.py & Correctly identifies presence/absence of MX and SPF records & Pass \\
blacklist\_checker.py & Returns match for known malicious domain; returns clean for safe domain & Pass \\
content\_checker.py & Detects urgency phrases, trust pages, and sensitive form fields correctly & Pass \\
behavior\_checker.py & Captures redirect chain and external domains; falls back gracefully without Playwright & Pass \\
privacy\_checker.py & Identifies tracking cookies, fingerprinting API patterns, and third-party scripts & Pass \\
tracker\_checker.py & Matches known tracker domains against local database and returns tracker name & Pass \\
scorer.py & Produces correct scores, labels, compound penalties, and blacklist override & Pass \\
\bottomrule
\end{tabular}
\caption{Unit Testing Results for ProblySUS Modules}
\end{table}

\section{Integration Testing}

Following unit testing, integration testing assessed the full analysis pipeline end-to-end — verifying that all modules interacted correctly, results were properly passed between components, and the final JSON response was correctly consumed by both the web frontend and the browser extension.

\textbf{Scenario 1 — Known Malicious URL (Blacklist Hit):}  
A domain present in the URLhaus feed was submitted. The blacklist checker correctly identified it, the scorer applied the hard override, and the frontend rendered a score of 100 with the \textit{Fraudulent} label.

\textbf{Scenario 2 — Newly Registered Phishing-Pattern Domain:}  
A URL using an abused TLD (.xyz), multiple hyphens, and registered within 7 days was submitted. The pipeline accumulated penalties across domain age, URL patterns, and DNS configuration, producing a Suspicious/Fraudulent verdict.

\textbf{Scenario 3 — Legitimate Established Website:}  
A well-known website was submitted. The system confirmed HTTPS validity, established domain age, and clean network behavior, producing a Safe score.

\textbf{Scenario 4 — JavaScript Redirect Site:}  
A site performing JavaScript-based redirects was submitted. The Playwright behavior module successfully detected redirect chains and captured external domains.

\textbf{Scenario 5 — Fast Mode vs Full Mode:}  
The same URL was tested with and without \texttt{?fast=true}. Fast mode produced a result in approximately 4 seconds while full mode produced a more detailed result in around 10 seconds.

All integration tests yielded stable results across repeated runs.

\section{Performance Testing}

Performance testing evaluated system responsiveness, concurrency handling, and reliability under realistic scanning conditions.

\begin{table}[h]
\centering
\begin{tabular}{|ll|}
\toprule
Test Scenario & Result \\
\midrule
Fast Mode scan (Playwright skipped) & Completed in 3--5 seconds \\
Full scan (Playwright included) & Completed in 8--12 seconds \\
Cached URL re-scan & Returned instantly ($<$50 ms) \\
Simultaneous requests & Handled concurrently without crashes \\
URLhaus blacklist auto-refresh & Completed in background without blocking requests \\
Behavior module timeout & Graceful fallback after 6 seconds \\
Average API response time (fast mode) & 350--700 ms excluding Playwright \\
\bottomrule
\end{tabular}
\caption{Performance Testing Results}
\end{table}

These results confirm the system's readiness for real-world usage.

\section{Usability Testing}

Usability testing was performed with a group of students and general internet users to assess the ease of use and interpretability of results.

\begin{itemize}
\item Approximately 90\% of testers found the risk label and color-coded score immediately understandable.
\item Around 85\% considered the reasons list helpful in understanding why a site received its score.
\item 80\% successfully used the browser extension to scan a site directly from their browser tab on the first attempt.
\item 100\% of participants expressed interest in using the tool regularly when encountering unfamiliar URLs.
\item Testers identified the Privacy Panel and Tracker Panel as particularly informative features.
\end{itemize}

\section{Bug Fixes and Improvements}

During testing, several bugs and performance issues were identified and resolved.

\begin{table}[H]
\centering
\begin{tabular}{|ll|}
\toprule
Issue Identified & Action Taken \\
\midrule
URLhaus feed ZIP caused CSV reader failure & Added zipfile extraction before parsing \\
Behavior module timeout on blocked sites & Added 6-second hard timeout and fallback \\
False positives on obfuscated JS & Made penalty context-aware for older domains \\
Redirect penalty too aggressive & Increased redirect threshold to 7+ \\
Extension popup crash when backend offline & Added try/catch and user-friendly error message \\
Multiple HTML fetches by modules & Refactored to fetch HTML once in app.py \\
Privacy exposure bar not animating & Reset animation state before rendering \\
\bottomrule
\end{tabular}
\caption{Bug Fixes and Enhancements}
\end{table}

These improvements significantly enhanced system reliability and user experience across multiple testing iterations.